\documentclass[a4paper,10pt]{article}

\usepackage[utf8]{inputenc}
\usepackage{t1enc}
\usepackage[spanish]{babel}
\usepackage[pdftex,usenames,dvipsnames]{color}
\usepackage[pdftex]{graphicx}
\usepackage{amsmath}
\usepackage{amsfonts}
\usepackage{amssymb}
\usepackage{listings}
\lstset{language=C}
\lstset{showstringspaces=false}
\lstset{basicstyle=\ttfamily,}

\begin{document}


\renewcommand{\lstlistingname}{C\'odigo Fuente}
\lstloadlanguages{Octave} 
\lstdefinelanguage{MyPseudoCode}[]{Octave}{
	deletekeywords={beta,det},
	morekeywords={repmat}
} 
\lstset{
	language=MyPseudoCode,
	stringstyle=\ttfamily,
	showstringspaces = false,
	basicstyle=\footnotesize\ttfamily,
	commentstyle=\color{gray},
	keywordstyle=\bfseries,
	numbers=left,
	numberstyle=\ttfamily\footnotesize,
	stepnumber=1,                   
	framexleftmargin=0.20cm,
	numbersep=0.37cm,              
	backgroundcolor=\color{white},
	showspaces=false,
	showtabs=false,
	frame=l,
	tabsize=4,
	captionpos=b,               
	breaklines=true,             
	breakatwhitespace=false,      
	mathescape=true
}
\begin{titlepage}
        \thispagestyle{empty}
        \begin{center}
                \includegraphics{./images/itba.jpg}
                \vfill
                \Huge{Sistemas Operativos}\\
                \vspace{1cm}
                \huge{Trabajo Práctico Especial Nº1}\\
        \end{center}
        \vspace{2cm}
        \large{
                \begin{tabular}{lcrc}
                        \textbf{Alvaro Crespo} & & 50758 & \ \ \texttt{acrespo@alu.itba.edu.ar}\\
                        \textbf{Juan Pablo Civile} & & 50453 & \ \ \texttt{jcivile@alu.itba.edu.ar}\\
                        \textbf{Darío Susnisky} & & 50592 & \ \ \texttt{dsusnisk@alu.itba.edu.ar}\\
                        \\ 
                \end{tabular}
        }
        \vfill
        \flushright{12 de Septiembre del 2011}
\end{titlepage}

\setcounter{page}{1}

\tableofcontents
\newpage
\section{Introducción}
El objetivo de este trabajo es familiarizarse con el uso de sistemas cliente-servidor concurrentes, implementando el servidor mediante la creación de procesos hijos 
utilizando \textit{fork()} y mediante la creación de \textit{threads}. Al mismo tiempo, ejercitar el uso de los distintos tipos de primitivas de sincronización y 
comunicación de procesos (IPC) y manejar con autoridad el \textit{filesystem} de Linux desde el lado usuario.

\newpage
\section{Esquema de la aplicación}
Esquema onda el que hicimos con bruno sobre que cosas son procesos y que son threads, y que cosa se comunica con que. Aunque sea un paint.

\newpage
\section{Modelo OSI}
Separación en capas. Separación clara de responsabilidades. Mejora en la claridad del código.
Esquema gráfico.

\newpage
\section{Interfaz de comunicación(???)}
Presentar las distintas interfaces que tuvimos...
Creo que fueron tres en total... 
Criterios de elección Y de descarte...
CHAMPO esta es TU parte, te llama a gritosss jeje.

\newpage
\section{Problemas encontrados(???)}
problema con el formato de los archivos de configuración. Formato poco feliz para C. Debimos implementaron una estructura de datos especial. 
problema con las líneas en blanco que marcaban los límites. fscanf(``\%s'',str) se comia el \\n de la linea en blanco. 

\newpage
\section{Conclusiones}
ACUERDENSE DE COMO NOS LA HIZO COMER ROMANNNNN!!!PONGAMOS ALGO COPADO ACA!!!XD

\newpage     
\section{Referencias}

\begin{itemize}
  \item Material provisto por la cátedra.
  \item UNIX system programming. Second Edition. Keith Havilland, Dina Gray, Ben Salama.
  \item http://cplusplus.com/
  \item http://beej.us/guide/bgipc/output/html/multipage/unixsock.html
  \item https://computing.llnl.gov/tutorials/pthreads/
  \item https://computing.llnl.gov/tutorials/parallel\_comp/
  \item http://www.csc.villanova.edu/~mdamian/threads/posixsem.html
  \item http://www.cs.cf.ac.uk/Dave/C/node25.html
  \item http://www.users.pjwstk.edu.pl/~jms/qnx/help/watcom/clibref/mq\_overview.html
  \item http://linux.die.net/man/7/mq\_overview
\end{itemize}
   
\end{document}
